% Ngân hàng bài dummy (dữ liệu test): Hàm số bậc nhất (lớp 8).
\documentclass[12pt,a4paper]{article}
\usepackage[utf8]{vietnam}
\usepackage{amsmath,amssymb}
\newenvironment{problem}[1][]{\par\medskip\noindent\textbf{Bài #1.}\ }{\par}
\newenvironment{solution}[1][]{\par\noindent\textit{Lời giải #1.}\par}{\par}

\begin{document}

% ---------------------------------------------------------------- giá trị hàm số
\begin{problem}[8DS-HSBN-001]
% nguon: dummy (tự sinh để test)
% lop: 8
% kc_du_kien: 8DS-HSBN-02
% do_kho_du_kien: 1
Cho hàm số $y = f(x) = 3x - 2$. Tính $f(0)$, $f(2)$ và $f\left(-\dfrac{1}{3}\right)$.

\begin{solution}
Để tính giá trị của hàm số tại $x = x_0$, ta thay $x = x_0$ vào công thức $f(x) = 3x - 2$ rồi thực hiện phép tính.

Thay $x = 0$: $f(0) = 3 \cdot 0 - 2 = -2$.

Thay $x = 2$: $f(2) = 3 \cdot 2 - 2 = 6 - 2 = 4$.

Thay $x = -\dfrac{1}{3}$: $f\left(-\dfrac{1}{3}\right) = 3 \cdot \left(-\dfrac{1}{3}\right) - 2 = -1 - 2 = -3$.

Đáp số: $\boxed{f(0) = -2,\ f(2) = 4,\ f\left(-\dfrac{1}{3}\right) = -3}$
\end{solution}
\end{problem}

% ---------------------------------------------------------------- nhận biết hàm số bậc nhất, nhiều ý
\begin{problem}[8DS-HSBN-002]
% nguon: dummy (tự sinh để test)
% lop: 8
% kc_du_kien: 8DS-HSBN-02
% do_kho_du_kien: 2
a) Trong các hàm số sau, hàm số nào là hàm số bậc nhất? Với mỗi hàm số bậc nhất, hãy chỉ ra hệ số $a$ và $b$.
(1) $y = 5 - 2x$; (2) $y = 2x^2 + 1$; (3) $y = \dfrac{x}{3}$; (4) $y = 2(x - 1) - 2x + 3$.

b) Tìm $m$ để hàm số $y = (2m - 4)x + 1$ là hàm số bậc nhất.

\begin{solution}
a) Hàm số bậc nhất là hàm số cho bởi công thức $y = ax + b$, trong đó $a$, $b$ là các số cho trước và $a \neq 0$.

a) Hàm số (1): $y = 5 - 2x = -2x + 5$ có dạng $y = ax + b$ với $a = -2 \neq 0$, $b = 5$, nên là hàm số bậc nhất.

a) Hàm số (2): $y = 2x^2 + 1$ chứa $x^2$ nên không viết được dưới dạng $y = ax + b$; nó không là hàm số bậc nhất.

a) Hàm số (3): $y = \dfrac{x}{3} = \dfrac{1}{3}x + 0$ có dạng $y = ax + b$ với $a = \dfrac{1}{3} \neq 0$, $b = 0$, nên là hàm số bậc nhất.

a) Hàm số (4): rút gọn vế phải theo quy tắc bỏ ngoặc, $y = 2x - 2 - 2x + 3 = 0x + 1$; hệ số của $x$ bằng $0$ nên không là hàm số bậc nhất (đây là hàm hằng $y = 1$).

b) Hàm số $y = (2m - 4)x + 1$ có dạng $y = ax + b$ với $a = 2m - 4$, $b = 1$; nó là hàm số bậc nhất khi và chỉ khi $a = 2m - 4 \neq 0$.

b) Giải điều kiện $2m - 4 \neq 0$: chuyển vế được $2m \neq 4$, chia hai vế cho $2$ được $m \neq 2$.

Đáp số: a) $\boxed{\text{(1): } a = -2,\ b = 5;\ \text{(3): } a = \dfrac{1}{3},\ b = 0}$; b) $\boxed{m \neq 2}$
\end{solution}
\end{problem}

% ---------------------------------------------------------------- điểm thuộc đồ thị, giao với hai trục, diện tích
\begin{problem}[8DS-HSBN-003]
% nguon: dummy (tự sinh để test)
% lop: 8
% kc_du_kien: 8DS-HSBN-01, 8DS-HSBN-03
% do_kho_du_kien: 2
Cho hàm số $y = -2x + 3$ có đồ thị là đường thẳng $d$.

a) Các điểm $M(1; 1)$ và $N(2; 1)$ có thuộc $d$ hay không?

b) Tìm toạ độ giao điểm $A$ của $d$ với trục tung và giao điểm $B$ của $d$ với trục hoành. Từ đó mô tả cách vẽ $d$.

c) Tính diện tích tam giác $OAB$ ($O$ là gốc toạ độ, đơn vị đo trên các trục là xentimét).

\begin{solution}
a) Điểm $(x_0; y_0)$ thuộc $d$ khi và chỉ khi toạ độ của nó thoả mãn công thức hàm số, tức là $y_0 = -2x_0 + 3$.

a) Với $M(1; 1)$: $-2 \cdot 1 + 3 = 1$, đúng bằng tung độ của $M$, nên $M$ thuộc $d$.

a) Với $N(2; 1)$: $-2 \cdot 2 + 3 = -1 \neq 1$, nên $N$ không thuộc $d$.

b) Giao điểm với trục tung có hoành độ bằng $0$; thay $x = 0$ được $y = -2 \cdot 0 + 3 = 3$, nên $A(0; 3)$.

b) Giao điểm với trục hoành có tung độ bằng $0$; cho $y = 0$ được $-2x + 3 = 0$, suy ra $x = \dfrac{3}{2}$, nên $B\left(\dfrac{3}{2}; 0\right)$.

b) Đồ thị của hàm số bậc nhất là một đường thẳng, nên $d$ là đường thẳng đi qua hai điểm $A(0; 3)$ nằm trên tia $Oy$ và $B\left(\dfrac{3}{2}; 0\right)$ nằm trên tia $Ox$. Để vẽ $d$, ta biểu diễn hai điểm $A$, $B$ trên mặt phẳng toạ độ rồi kẻ đường thẳng $AB$.

c) Vì hai trục toạ độ vuông góc nên tam giác $OAB$ vuông tại $O$, với hai cạnh góc vuông $OA = 3$ cm và $OB = \dfrac{3}{2}$ cm (độ dài bằng giá trị tuyệt đối của toạ độ khác $0$ của $A$, $B$).

c) Diện tích tam giác vuông bằng nửa tích hai cạnh góc vuông: $S = \dfrac{1}{2} \cdot OA \cdot OB = \dfrac{1}{2} \cdot 3 \cdot \dfrac{3}{2} = \dfrac{9}{4}$ (cm$^2$).

Đáp số: a) $\boxed{M \in d,\ N \notin d}$; b) $\boxed{A(0; 3),\ B\left(\dfrac{3}{2}; 0\right)}$; c) $\boxed{S = \dfrac{9}{4} \text{ cm}^2}$
\end{solution}
\end{problem}

% ---------------------------------------------------------------- vị trí hai đường thẳng theo tham số
\begin{problem}[8DS-HSBN-004]
% nguon: dummy (tự sinh để test)
% lop: 8
% kc_du_kien: 8DS-HSBN-04
% do_kho_du_kien: 3
Cho hai đường thẳng $d_1: y = (m + 1)x + 3$ (với $m \neq -1$) và $d_2: y = 2x + m$. Tìm $m$ để:

a) $d_1$ song song với $d_2$;

b) $d_1$ cắt $d_2$;

c) $d_1$ trùng với $d_2$.

\begin{solution}
a) Hai đường thẳng $y = ax + b$ và $y = a'x + b'$ (với $a \neq 0$, $a' \neq 0$) song song khi và chỉ khi $a = a'$ và $b \neq b'$. Ở đây $a = m + 1$, $b = 3$, $a' = 2$, $b' = m$.

a) Điều kiện $a = a'$: $m + 1 = 2$, suy ra $m = 1$ (thoả mãn $m \neq -1$).

a) Với $m = 1$ thì $b' = 1 \neq 3 = b$, nên điều kiện $b \neq b'$ cũng thoả mãn. Vậy $d_1 \parallel d_2$ khi và chỉ khi $m = 1$.

b) Hai đường thẳng cắt nhau khi và chỉ khi hai hệ số góc khác nhau: $m + 1 \neq 2$, tức là $m \neq 1$. Kết hợp với điều kiện $m \neq -1$ của đề bài.

c) Hai đường thẳng trùng nhau khi và chỉ khi $a = a'$ và $b = b'$, tức là đồng thời $m + 1 = 2$ và $3 = m$.

c) Từ $m + 1 = 2$ có $m = 1$, mâu thuẫn với $m = 3$, nên không có giá trị nào của $m$ để $d_1$ trùng với $d_2$.

Đáp số: a) $\boxed{m = 1}$; b) $\boxed{m \neq 1 \text{ và } m \neq -1}$; c) $\boxed{\text{Không có giá trị nào của } m}$
\end{solution}
\end{problem}

% ---------------------------------------------------------------- xác định hàm số
\begin{problem}[8DS-HSBN-005]
% nguon: dummy (tự sinh để test)
% lop: 8
% kc_du_kien: 8DS-HSBN-05, 8DS-HSBN-04, 8DS-HSBN-01
% do_kho_du_kien: 2
Xác định hàm số bậc nhất $y = ax + b$, biết đồ thị của nó song song với đường thẳng $y = -3x + 5$ và cắt trục hoành tại điểm có hoành độ bằng $2$.

\begin{solution}
Đồ thị hàm số $y = ax + b$ song song với đường thẳng $y = -3x + 5$ nên hai hệ số góc bằng nhau và hai tung độ gốc khác nhau: $a = -3$ và $b \neq 5$.

Đồ thị cắt trục hoành tại điểm có hoành độ bằng $2$, nghĩa là đồ thị đi qua điểm $(2; 0)$; do đó toạ độ điểm này thoả mãn công thức hàm số: $0 = a \cdot 2 + b$.

Thay $a = -3$: $0 = -6 + b$, suy ra $b = 6$.

Đối chiếu điều kiện: $b = 6 \neq 5$ (thoả mãn), và $a = -3 \neq 0$ nên hàm số tìm được là hàm số bậc nhất.

Đáp số: $\boxed{y = -3x + 6}$
\end{solution}
\end{problem}

% ---------------------------------------------------------------- bài toán thực tế, giao điểm hai đường thẳng
\begin{problem}[8DS-HSBN-006]
% nguon: dummy (tự sinh để test)
% lop: 8
% kc_du_kien: 8DS-HSBN-02, 8DS-HSBN-06
% do_kho_du_kien: 3
Hãng taxi A tính tiền như sau: $15$ nghìn đồng tiền mở cửa và $12$ nghìn đồng cho mỗi kilômét đi được. Hãng taxi B tính: $9$ nghìn đồng tiền mở cửa và $14$ nghìn đồng cho mỗi kilômét đi được.

a) Gọi $x$ (km) là quãng đường đi; $y_A$, $y_B$ (nghìn đồng) lần lượt là số tiền phải trả khi đi taxi hãng A, hãng B. Viết công thức tính $y_A$, $y_B$ theo $x$. Các hàm số này có phải là hàm số bậc nhất không?

b) Với quãng đường bao nhiêu kilômét thì số tiền phải trả cho hai hãng bằng nhau?

c) Một người cần đi quãng đường $10$ km. Người đó nên chọn hãng nào để tốn ít tiền hơn, và tiết kiệm được bao nhiêu?

\begin{solution}
a) Số tiền phải trả bằng tiền mở cửa cộng với tiền theo số kilômét, nên $y_A = 12x + 15$ và $y_B = 14x + 9$ (với $x > 0$).

a) Cả hai công thức đều có dạng $y = ax + b$ với hệ số của $x$ lần lượt là $12 \neq 0$ và $14 \neq 0$, nên đều là hàm số bậc nhất.

b) Số tiền hai hãng bằng nhau khi $y_A = y_B$, tức là $12x + 15 = 14x + 9$.

b) Chuyển vế: $15 - 9 = 14x - 12x$, nên $2x = 6$, suy ra $x = 3$ (thoả mãn $x > 0$).

b) Khi đó $y_A = y_B = 12 \cdot 3 + 15 = 51$ (nghìn đồng). Trên mặt phẳng toạ độ, $(3; 51)$ chính là giao điểm của hai đường thẳng $y = 12x + 15$ và $y = 14x + 9$.

c) Thay $x = 10$: $y_A = 12 \cdot 10 + 15 = 135$ (nghìn đồng) và $y_B = 14 \cdot 10 + 9 = 149$ (nghìn đồng).

c) Vì $135 < 149$ nên người đó nên chọn hãng A, tiết kiệm được $149 - 135 = 14$ (nghìn đồng).

Đáp số: a) $\boxed{y_A = 12x + 15,\ y_B = 14x + 9 \text{ (đều là hàm số bậc nhất)}}$; b) $\boxed{3 \text{ km}}$; c) $\boxed{\text{Chọn hãng A, tiết kiệm } 14 \text{ nghìn đồng}}$
\end{solution}
\end{problem}

% ---------------------------------------------------------------- giao điểm hai đường thẳng nằm trên trục, tham số
\begin{problem}[8DS-HSBN-007]
% nguon: dummy (tự sinh để test)
% lop: 8
% kc_du_kien: 8DS-HSBN-06, 8DS-HSBN-03
% do_kho_du_kien: 4
Cho hai đường thẳng $d_1: y = 2x + m - 1$ và $d_2: y = -x + 2m + 2$ (với $m$ là tham số).

a) Chứng minh rằng với mọi giá trị của $m$, $d_1$ và $d_2$ luôn cắt nhau. Tìm $m$ để $d_1$ và $d_2$ cắt nhau tại một điểm nằm trên trục tung.

b) Tìm $m$ để $d_1$ và $d_2$ cắt nhau tại một điểm nằm trên trục hoành.

\begin{solution}
a) Hệ số góc của $d_1$ là $2$, của $d_2$ là $-1$; vì $2 \neq -1$ nên $d_1$ và $d_2$ luôn cắt nhau tại đúng một điểm với mọi $m$.

a) Hoành độ giao điểm là nghiệm của phương trình hoành độ giao điểm $2x + m - 1 = -x + 2m + 2$; chuyển vế được $3x = m + 3$, suy ra $x = \dfrac{m + 3}{3}$.

a) Giao điểm nằm trên trục tung khi và chỉ khi hoành độ của nó bằng $0$: $\dfrac{m + 3}{3} = 0$, suy ra $m = -3$.

a) Thử lại: với $m = -3$ thì $d_1: y = 2x - 4$ và $d_2: y = -x - 4$ cùng đi qua điểm $(0; -4)$ trên trục tung (đúng).

b) Tung độ giao điểm: thay $x = \dfrac{m + 3}{3}$ vào phương trình của $d_1$ được $y = 2 \cdot \dfrac{m + 3}{3} + m - 1 = \dfrac{2m + 6 + 3m - 3}{3} = \dfrac{5m + 3}{3}$.

b) Giao điểm nằm trên trục hoành khi và chỉ khi tung độ của nó bằng $0$: $\dfrac{5m + 3}{3} = 0$, tức là $5m + 3 = 0$, suy ra $m = -\dfrac{3}{5}$.

b) Thử lại: với $m = -\dfrac{3}{5}$ thì $m + 3 = \dfrac{12}{5}$, hoành độ giao điểm $x = \dfrac{12}{5} : 3 = \dfrac{4}{5}$, và tung độ $y = 2 \cdot \dfrac{4}{5} - \dfrac{3}{5} - 1 = \dfrac{8 - 3 - 5}{5} = 0$ (đúng).

Đáp số: a) $\boxed{m = -3}$; b) $\boxed{m = -\dfrac{3}{5}}$
\end{solution}
\end{problem}

% ---------------------------------------------------------------- ba đường thẳng đồng quy
\begin{problem}[8DS-HSBN-008]
% nguon: dummy (tự sinh để test)
% lop: 8
% kc_du_kien: 8DS-HSBN-06, 8DS-HSBN-01
% do_kho_du_kien: 3
Cho ba đường thẳng $d_1: y = 2x - 1$, $d_2: y = -x + 5$ và $d_3: y = (m - 1)x + m$ (với $m \neq 1$).

a) Tìm toạ độ giao điểm $I$ của $d_1$ và $d_2$.

b) Tìm $m$ để ba đường thẳng $d_1$, $d_2$, $d_3$ đồng quy (cùng đi qua một điểm).

\begin{solution}
a) Vì hai hệ số góc $2 \neq -1$ nên $d_1$ và $d_2$ cắt nhau; hoành độ giao điểm là nghiệm của phương trình hoành độ giao điểm $2x - 1 = -x + 5$.

a) Chuyển vế: $3x = 6$, suy ra $x = 2$; thay vào phương trình của $d_1$: $y = 2 \cdot 2 - 1 = 3$. Kiểm tra trên $d_2$: $-2 + 5 = 3$ (đúng). Vậy $I(2; 3)$.

b) Vì $d_1$ và $d_2$ chỉ có một điểm chung là $I$, nên ba đường thẳng đồng quy khi và chỉ khi $d_3$ đi qua $I(2; 3)$.

b) $d_3$ đi qua $I(2; 3)$ khi và chỉ khi toạ độ của $I$ thoả mãn phương trình của $d_3$: $3 = (m - 1) \cdot 2 + m$.

b) Giải phương trình: $3 = 2m - 2 + m$, nên $3m = 5$, suy ra $m = \dfrac{5}{3}$ (thoả mãn $m \neq 1$).

b) Thử lại: với $m = \dfrac{5}{3}$ thì $d_3: y = \dfrac{2}{3}x + \dfrac{5}{3}$ và $\dfrac{2}{3} \cdot 2 + \dfrac{5}{3} = 3$, nên $d_3$ đi qua $I$; hệ số góc $\dfrac{2}{3}$ khác $2$ và $-1$ nên ba đường thẳng phân biệt.

Đáp số: a) $\boxed{I(2; 3)}$; b) $\boxed{m = \dfrac{5}{3}}$
\end{solution}
\end{problem}

% ---------------------------------------------------------------- điểm cố định, có cách giải khác
\begin{problem}[8DS-HSBN-009]
% nguon: dummy (tự sinh để test)
% lop: 8
% kc_du_kien: 8DS-HSBN-01, 8DS-HSBN-06
% do_kho_du_kien: 4
Cho hàm số $y = (m - 1)x + 2m - 3$ (với $m \neq 1$) có đồ thị là đường thẳng $d_m$. Chứng minh rằng khi $m$ thay đổi, đường thẳng $d_m$ luôn đi qua một điểm cố định. Tìm toạ độ điểm cố định đó.

\begin{solution}[cách 1: tách theo tham số]
Điểm $I(x_0; y_0)$ là điểm cố định mà $d_m$ luôn đi qua khi và chỉ khi $y_0 = (m - 1)x_0 + 2m - 3$ đúng với mọi $m \neq 1$ (điều kiện điểm thuộc đồ thị).

Bỏ ngoặc và nhóm các hạng tử chứa $m$: $y_0 = mx_0 - x_0 + 2m - 3$, tức là $m(x_0 + 2) - (x_0 + y_0 + 3) = 0$.

Nếu $x_0 + 2 \neq 0$ thì đẳng thức trên chỉ đúng với đúng một giá trị của $m$; vì vậy nó đúng với mọi $m \neq 1$ khi và chỉ khi $x_0 + 2 = 0$ và $x_0 + y_0 + 3 = 0$.

Từ $x_0 + 2 = 0$ có $x_0 = -2$; thay vào $x_0 + y_0 + 3 = 0$ được $y_0 = -1$.

Thử lại: với mọi $m \neq 1$, $(m - 1) \cdot (-2) + 2m - 3 = -2m + 2 + 2m - 3 = -1$, nên điểm $I(-2; -1)$ luôn thuộc $d_m$.

Đáp số: $\boxed{I(-2; -1)}$
\end{solution}

\begin{solution}[cách 2: giao điểm của hai đường thẳng cụ thể]
Nếu $d_m$ luôn đi qua điểm cố định $I$ thì $I$ thuộc cả $d_0$ (ứng với $m = 0$): $y = -x - 3$ và $d_2$ (ứng với $m = 2$): $y = x + 1$.

Hai đường thẳng $d_0$, $d_2$ có hệ số góc $-1 \neq 1$ nên cắt nhau tại đúng một điểm; hoành độ giao điểm là nghiệm của $-x - 3 = x + 1$, tức là $2x = -4$, $x = -2$; khi đó $y = -2 + 1 = -1$.

Vậy nếu có điểm cố định thì điểm đó chỉ có thể là $I(-2; -1)$.

Kiểm tra $I$ thuộc mọi $d_m$: $(m - 1) \cdot (-2) + 2m - 3 = -2m + 2 + 2m - 3 = -1$ với mọi $m \neq 1$, nên $d_m$ luôn đi qua $I(-2; -1)$.

Đáp số: $\boxed{I(-2; -1)}$
\end{solution}
\end{problem}

% ---------------------------------------------------------------- song song, tam giác vuông cân với hai trục
\begin{problem}[8DS-HSBN-010]
% nguon: dummy (tự sinh để test)
% lop: 8
% kc_du_kien: 8DS-HSBN-04, 8DS-HSBN-03
% do_kho_du_kien: 5
Cho đường thẳng $d: y = (m - 3)x + m - 1$ (với $m \neq 3$).

a) Tìm $m$ để $d$ song song với đường thẳng $d': y = -2x + 5$.

b) Tìm $m$ để $d$ cắt trục tung tại $A$, cắt trục hoành tại $B$ sao cho tam giác $OAB$ vuông cân ($O$ là gốc toạ độ).

\begin{solution}
a) $d \parallel d'$ khi và chỉ khi hai hệ số góc bằng nhau và hai tung độ gốc khác nhau: $m - 3 = -2$ và $m - 1 \neq 5$.

a) Từ $m - 3 = -2$ có $m = 1$; khi đó $m - 1 = 0 \neq 5$ (thoả mãn) và $m \neq 3$. Vậy $m = 1$, lúc đó $d: y = -2x$.

b) Giao điểm với trục tung: cho $x = 0$ được $y = m - 1$, nên $A(0; m - 1)$.

b) Để có tam giác $OAB$ thì $A$ phải khác $O$, tức là $m - 1 \neq 0$, hay $m \neq 1$. (Khi $m = 1$, $d$ đi qua gốc toạ độ nên $A$, $B$ đều trùng $O$.)

b) Giao điểm với trục hoành: cho $y = 0$ được $(m - 3)x + m - 1 = 0$; vì $m \neq 3$ nên $x = \dfrac{1 - m}{m - 3}$, do đó $B\left(\dfrac{1 - m}{m - 3}; 0\right)$.

b) Vì $A$ nằm trên trục tung, $B$ nằm trên trục hoành và hai trục vuông góc nên tam giác $OAB$ luôn vuông tại $O$; do đó nó vuông cân khi và chỉ khi $OA = OB$.

b) Ta có $OA = |m - 1|$ và $OB = \left|\dfrac{1 - m}{m - 3}\right| = \dfrac{|m - 1|}{|m - 3|}$. Điều kiện $OA = OB$ trở thành $|m - 1| = \dfrac{|m - 1|}{|m - 3|}$.

b) Vì $m \neq 1$ nên $|m - 1| > 0$; chia hai vế cho $|m - 1|$ rồi nhân hai vế với $|m - 3|$ được $|m - 3| = 1$, suy ra $m - 3 = 1$ hoặc $m - 3 = -1$, tức là $m = 4$ hoặc $m = 2$.

b) Đối chiếu điều kiện $m \neq 1$, $m \neq 3$: cả hai giá trị đều thoả mãn. Thử lại: $m = 4$ cho $d: y = x + 3$, $A(0; 3)$, $B(-3; 0)$, $OA = OB = 3$; $m = 2$ cho $d: y = -x + 1$, $A(0; 1)$, $B(1; 0)$, $OA = OB = 1$ (đúng).

Đáp số: a) $\boxed{m = 1}$; b) $\boxed{m = 2 \text{ hoặc } m = 4}$
\end{solution}
\end{problem}

% ---------------------------------------------------------------- hệ số góc, góc tạo với trục Ox
\begin{problem}[8DS-HSBN-011]
% nguon: dummy (tự sinh để test)
% lop: 8
% kc_du_kien: 8DS-HSBN-04
% do_kho_du_kien: 1
Cho các hàm số bậc nhất: (1) $y = 4x - 7$; (2) $y = 3 - \dfrac{2}{5}x$; (3) $y = \dfrac{x + 6}{2}$; (4) $y = -(x + 1)$. Với mỗi hàm số, hãy chỉ ra hệ số góc của đồ thị và cho biết góc tạo bởi đồ thị đó với trục $Ox$ là góc nhọn hay góc tù.

\begin{solution}
Đồ thị của hàm số bậc nhất $y = ax + b$ ($a \neq 0$) là đường thẳng có hệ số góc $a$; góc tạo bởi đường thẳng đó và trục $Ox$ là góc nhọn khi $a > 0$ và là góc tù khi $a < 0$. Vì vậy ta viết mỗi hàm số về dạng $y = ax + b$ để đọc ra $a$.

Hàm số (1): $y = 4x - 7$ có $a = 4 > 0$, nên hệ số góc là $4$ và góc tạo với trục $Ox$ là góc nhọn.

Hàm số (2): đổi chỗ các hạng tử, $y = 3 - \dfrac{2}{5}x = -\dfrac{2}{5}x + 3$ có $a = -\dfrac{2}{5} < 0$, nên hệ số góc là $-\dfrac{2}{5}$ và góc tạo với trục $Ox$ là góc tù.

Hàm số (3): chia từng hạng tử của tử cho $2$, $y = \dfrac{x + 6}{2} = \dfrac{1}{2}x + 3$ có $a = \dfrac{1}{2} > 0$, nên hệ số góc là $\dfrac{1}{2}$ và góc tạo với trục $Ox$ là góc nhọn.

Hàm số (4): bỏ ngoặc có dấu trừ đằng trước, $y = -(x + 1) = -x - 1$ có $a = -1 < 0$, nên hệ số góc là $-1$ và góc tạo với trục $Ox$ là góc tù.

Đáp số: $\boxed{(1)\ a = 4 \text{ (góc nhọn)};\ (2)\ a = -\dfrac{2}{5} \text{ (góc tù)};\ (3)\ a = \dfrac{1}{2} \text{ (góc nhọn)};\ (4)\ a = -1 \text{ (góc tù)}}$
\end{solution}
\end{problem}

% ---------------------------------------------------------------- tìm tham số để đồ thị đi qua điểm; điểm có tung độ gấp đôi hoành độ
\begin{problem}[8DS-HSBN-012]
% nguon: dummy (tự sinh để test)
% lop: 8
% kc_du_kien: 8DS-HSBN-01, 8DS-HSBN-02
% do_kho_du_kien: 2
Cho hàm số $y = (2m - 1)x + m + 2$, với $m$ là tham số.

a) Tìm $m$ để hàm số đã cho là hàm số bậc nhất và đồ thị của nó đi qua điểm $A(-1; 4)$.

b) Với giá trị $m$ tìm được ở câu a), tìm toạ độ điểm $M$ thuộc đồ thị có tung độ bằng hai lần hoành độ.

\begin{solution}
a) Hàm số có dạng $y = ax + b$ với $a = 2m - 1$; nó là hàm số bậc nhất khi và chỉ khi $2m - 1 \neq 0$, tức là $m \neq \dfrac{1}{2}$.

a) Đồ thị đi qua $A(-1; 4)$ khi và chỉ khi toạ độ của $A$ thoả mãn công thức hàm số: $4 = (2m - 1) \cdot (-1) + m + 2$.

a) Bỏ ngoặc và thu gọn vế phải: $4 = -2m + 1 + m + 2 = -m + 3$, suy ra $m = -1$ (thoả mãn $m \neq \dfrac{1}{2}$). Khi đó hàm số là $y = -3x + 1$.

b) Điểm có tung độ bằng hai lần hoành độ có dạng $M(x_0; 2x_0)$. Điểm $M$ thuộc đồ thị $y = -3x + 1$ khi và chỉ khi $2x_0 = -3x_0 + 1$.

b) Chuyển vế: $5x_0 = 1$, suy ra $x_0 = \dfrac{1}{5}$ và tung độ $2x_0 = \dfrac{2}{5}$. Thử lại: $-3 \cdot \dfrac{1}{5} + 1 = \dfrac{2}{5}$ (đúng), nên $M\left(\dfrac{1}{5}; \dfrac{2}{5}\right)$.

Đáp số: a) $\boxed{m = -1}$; b) $\boxed{M\left(\dfrac{1}{5}; \dfrac{2}{5}\right)}$
\end{solution}
\end{problem}

% ---------------------------------------------------------------- xác định hàm số biết tung độ gốc và một điểm
\begin{problem}[8DS-HSBN-013]
% nguon: dummy (tự sinh để test)
% lop: 8
% kc_du_kien: 8DS-HSBN-05, 8DS-HSBN-03
% do_kho_du_kien: 2
Xác định hàm số bậc nhất $y = ax + b$, biết đồ thị của nó cắt trục tung tại điểm có tung độ bằng $-2$ và đi qua điểm $A(3; 4)$. Với hàm số tìm được, tìm toạ độ giao điểm của đồ thị với trục hoành.

\begin{solution}
Thay $x = 0$ vào $y = ax + b$ được $y = b$, nên đồ thị cắt trục tung tại điểm $(0; b)$. Theo đề bài, tung độ của điểm này bằng $-2$, do đó $b = -2$.

Đồ thị đi qua $A(3; 4)$ nên toạ độ của $A$ thoả mãn công thức hàm số: $4 = a \cdot 3 + b = 3a - 2$.

Chuyển vế: $3a = 6$, suy ra $a = 2$; vì $a = 2 \neq 0$ nên hàm số cần tìm là hàm số bậc nhất $y = 2x - 2$.

Giao điểm của đồ thị với trục hoành có tung độ bằng $0$; cho $y = 0$ được $2x - 2 = 0$, suy ra $x = 1$. Vậy giao điểm là $(1; 0)$.

Đáp số: $\boxed{y = 2x - 2;\ \text{giao điểm với trục hoành là } (1; 0)}$
\end{solution}
\end{problem}

% ---------------------------------------------------------------- đường thẳng qua hai điểm, ba điểm thẳng hàng
\begin{problem}[8DS-HSBN-014]
% nguon: dummy (tự sinh để test)
% lop: 8
% kc_du_kien: 8DS-HSBN-05, 8DS-HSBN-01
% do_kho_du_kien: 3
Cho ba điểm $A(1; 3)$, $B(-2; -3)$ và $C(3; 7)$.

a) Xác định hàm số bậc nhất $y = ax + b$ có đồ thị là đường thẳng đi qua hai điểm $A$ và $B$.

b) Chứng minh rằng ba điểm $A$, $B$, $C$ thẳng hàng.

c) Tìm $k$ để điểm $D(2k; k - 3)$ cũng thuộc đường thẳng $AB$.

\begin{solution}
a) Đường thẳng $y = ax + b$ đi qua $A(1; 3)$ nên $3 = a + b$ $(1)$; đi qua $B(-2; -3)$ nên $-3 = -2a + b$ $(2)$.

a) Từ $(1)$ suy ra $b = 3 - a$; thay vào $(2)$ được $-3 = -2a + 3 - a$, tức là $-3 = -3a + 3$.

a) Giải phương trình bậc nhất ẩn $a$: $3a = 6$, nên $a = 2$; suy ra $b = 3 - 2 = 1$. Vì $a = 2 \neq 0$ nên hàm số cần tìm là $y = 2x + 1$.

b) Ba điểm $A$, $B$, $C$ thẳng hàng khi và chỉ khi $C$ thuộc đường thẳng $AB$. Thay $x = 3$ vào $y = 2x + 1$ được $y = 7$, đúng bằng tung độ của $C$, nên $C$ thuộc đường thẳng $AB$; vậy $A$, $B$, $C$ thẳng hàng.

c) Điểm $D(2k; k - 3)$ thuộc đường thẳng $AB$ khi và chỉ khi toạ độ của $D$ thoả mãn $y = 2x + 1$: $k - 3 = 2 \cdot 2k + 1$.

c) Giải phương trình: $k - 3 = 4k + 1$, chuyển vế được $-3k = 4$, suy ra $k = -\dfrac{4}{3}$. Thử lại: $D\left(-\dfrac{8}{3}; -\dfrac{13}{3}\right)$ và $2 \cdot \left(-\dfrac{8}{3}\right) + 1 = -\dfrac{13}{3}$ (đúng).

Đáp số: a) $\boxed{y = 2x + 1}$; b) $\boxed{C \text{ thuộc đường thẳng } AB \text{ nên } A, B, C \text{ thẳng hàng}}$; c) $\boxed{k = -\dfrac{4}{3}}$
\end{solution}
\end{problem}

% ---------------------------------------------------------------- bài toán thực tế: bơm nước, xác định hàm số từ hai dữ kiện, hai cách
\begin{problem}[8DS-HSBN-015]
% nguon: dummy (tự sinh để test)
% lop: 8
% kc_du_kien: 8DS-HSBN-05, 8DS-HSBN-02
% do_kho_du_kien: 3
Người ta dùng máy bơm nước vào một bể. Lượng nước $y$ (lít) có trong bể sau $x$ giờ bơm ($x \geq 0$) là một hàm số bậc nhất của $x$: $y = ax + b$. Sau $2$ giờ bơm, trong bể có $1\,800$ lít nước; sau $5$ giờ bơm, trong bể có $3\,600$ lít nước.

a) Xác định $a$ và $b$.

b) Lúc bắt đầu bơm, trong bể đã có bao nhiêu lít nước? Mỗi giờ máy bơm được bao nhiêu lít nước?

c) Bể chứa được tối đa $6\,000$ lít nước. Sau bao nhiêu giờ bơm thì bể đầy?

\begin{solution}[cách 1: lập phương trình]
a) Với $x = 2$ thì $y = 1\,800$, nên $1\,800 = 2a + b$ $(1)$; với $x = 5$ thì $y = 3\,600$, nên $3\,600 = 5a + b$ $(2)$.

a) Từ $(1)$ suy ra $b = 1\,800 - 2a$; thay vào $(2)$: $3\,600 = 5a + 1\,800 - 2a = 3a + 1\,800$, nên $3a = 1\,800$, $a = 600$; khi đó $b = 1\,800 - 2 \cdot 600 = 600$. Vì $a \neq 0$ nên $y = 600x + 600$ là hàm số bậc nhất.

b) Lúc bắt đầu bơm thì $x = 0$, lượng nước có trong bể là $y = 600 \cdot 0 + 600 = 600$ (lít).

b) Khi $x$ tăng thêm $1$ (thêm một giờ bơm) thì $y$ tăng thêm $600(x + 1) + 600 - (600x + 600) = 600$; vậy mỗi giờ máy bơm được $600$ lít nước.

c) Bể đầy khi $y = 6\,000$, tức là $600x + 600 = 6\,000$; chuyển vế được $600x = 5\,400$, suy ra $x = 9$ (thoả mãn $x \geq 0$).

Đáp số: a) $\boxed{a = 600,\ b = 600}$; b) $\boxed{600 \text{ lít};\ 600 \text{ lít mỗi giờ}}$; c) $\boxed{9 \text{ giờ}}$
\end{solution}

\begin{solution}[cách 2: dùng ý nghĩa của hệ số $a$]
a) Với $y = ax + b$, khi $x$ tăng thêm $1$ thì $y$ tăng thêm $a(x + 1) + b - (ax + b) = a$; vậy $a$ là số lít nước bơm được trong mỗi giờ.

a) Từ thời điểm $2$ giờ đến thời điểm $5$ giờ (trong $3$ giờ), lượng nước tăng $3\,600 - 1\,800 = 1\,800$ (lít); mỗi giờ tăng $a$ lít nên $3a = 1\,800$, suy ra $a = 600$.

a) Thay $x = 2$, $y = 1\,800$ vào $y = 600x + b$: $1\,800 = 1\,200 + b$, suy ra $b = 600$.

b) $b$ là giá trị của $y$ khi $x = 0$, nên lúc bắt đầu bơm trong bể có $600$ lít nước; theo trên, mỗi giờ máy bơm được $a = 600$ lít nước.

c) Kể từ lúc bắt đầu, cần bơm thêm $6\,000 - 600 = 5\,400$ (lít) thì bể đầy; mỗi giờ bơm được $600$ lít nên cần $5\,400 : 600 = 9$ (giờ).

Đáp số: a) $\boxed{a = 600,\ b = 600}$; b) $\boxed{600 \text{ lít};\ 600 \text{ lít mỗi giờ}}$; c) $\boxed{9 \text{ giờ}}$
\end{solution}
\end{problem}

% ---------------------------------------------------------------- tam giác tạo bởi hai đường thẳng và trục hoành
\begin{problem}[8DS-HSBN-016]
% nguon: dummy (tự sinh để test)
% lop: 8
% kc_du_kien: 8DS-HSBN-06, 8DS-HSBN-03
% do_kho_du_kien: 3
Cho hai đường thẳng $d_1: y = x + 2$ và $d_2: y = -2x + 8$.

a) Tìm toạ độ giao điểm $C$ của $d_1$ và $d_2$.

b) Gọi $A$, $B$ lần lượt là giao điểm của $d_1$, $d_2$ với trục hoành. Tìm toạ độ của $A$, $B$ và tính diện tích tam giác $ABC$ (đơn vị đo trên các trục là xentimét).

\begin{solution}
a) Hai hệ số góc $1 \neq -2$ nên $d_1$ và $d_2$ cắt nhau; hoành độ giao điểm là nghiệm của phương trình hoành độ giao điểm $x + 2 = -2x + 8$.

a) Chuyển vế: $3x = 6$, suy ra $x = 2$; thay vào phương trình của $d_1$: $y = 2 + 2 = 4$. Kiểm tra trên $d_2$: $-2 \cdot 2 + 8 = 4$ (đúng). Vậy $C(2; 4)$.

b) Giao điểm với trục hoành có tung độ bằng $0$. Với $d_1$: $x + 2 = 0$, suy ra $x = -2$, nên $A(-2; 0)$. Với $d_2$: $-2x + 8 = 0$, suy ra $x = 4$, nên $B(4; 0)$.

b) Hai điểm $A$, $B$ cùng nằm trên trục hoành nên độ dài $AB$ bằng hiệu hai hoành độ (số lớn trừ số nhỏ): $AB = 4 - (-2) = 6$ (cm).

b) Gọi $H(2; 0)$ là hình chiếu của $C$ trên trục hoành; $CH$ vuông góc với $AB$ nên $CH$ là đường cao của tam giác $ABC$ ứng với cạnh $AB$, và $CH = 4$ cm (bằng tung độ của $C$).

b) Diện tích tam giác bằng nửa tích một cạnh với đường cao ứng với cạnh đó: $S_{ABC} = \dfrac{1}{2} \cdot AB \cdot CH = \dfrac{1}{2} \cdot 6 \cdot 4 = 12$ (cm$^2$).

Đáp số: a) $\boxed{C(2; 4)}$; b) $\boxed{A(-2; 0),\ B(4; 0),\ S_{ABC} = 12 \text{ cm}^2}$
\end{solution}
\end{problem}

% ---------------------------------------------------------------- giao điểm có toạ độ nguyên, tham số nguyên
\begin{problem}[8DS-HSBN-017]
% nguon: dummy (tự sinh để test)
% lop: 8
% kc_du_kien: 8DS-HSBN-06, 8DS-HSBN-04
% do_kho_du_kien: 4
Cho hàm số bậc nhất $y = (m - 1)x + 2$ (với $m \neq 1$) có đồ thị là đường thẳng $d_1$ và hàm số $y = 3x - 1$ có đồ thị là đường thẳng $d_2$. Tìm tất cả các giá trị nguyên của $m$ để $d_1$ cắt $d_2$ tại một điểm có hoành độ và tung độ đều là số nguyên.

\begin{solution}
Hai đường thẳng cắt nhau khi và chỉ khi hai hệ số góc khác nhau: $m - 1 \neq 3$, tức là $m \neq 4$; kết hợp với điều kiện $m \neq 1$ của đề bài.

Hoành độ giao điểm là nghiệm của phương trình hoành độ giao điểm $(m - 1)x + 2 = 3x - 1$; chuyển vế được $(m - 1 - 3)x = -1 - 2$, tức là $(m - 4)x = -3$.

Vì $m \neq 4$ nên $x = \dfrac{-3}{m - 4}$, và tung độ giao điểm là $y = 3x - 1$.

Nếu $x$ là số nguyên thì $y = 3x - 1$ cũng là số nguyên, nên chỉ cần $x$ nguyên. Với $m$ nguyên, $m - 4$ là số nguyên khác $0$, nên $x = \dfrac{-3}{m - 4}$ là số nguyên khi và chỉ khi $m - 4$ là ước của $3$, tức là $m - 4 \in \{-3; -1; 1; 3\}$.

Suy ra $m \in \{1; 3; 5; 7\}$. Đối chiếu điều kiện $m \neq 1$ (để $y = (m - 1)x + 2$ là hàm số bậc nhất) và $m \neq 4$: loại $m = 1$, còn lại $m \in \{3; 5; 7\}$.

Thử lại: $m = 3$ cho $x = 3$, $y = 8$ (và $2 \cdot 3 + 2 = 8$); $m = 5$ cho $x = -3$, $y = -10$ (và $4 \cdot (-3) + 2 = -10$); $m = 7$ cho $x = -1$, $y = -4$ (và $6 \cdot (-1) + 2 = -4$). Các giao điểm đều có toạ độ nguyên.

Đáp số: $\boxed{m \in \{3; 5; 7\}}$
\end{solution}
\end{problem}

% ---------------------------------------------------------------- tìm hàm số bậc nhất từ đẳng thức hàm hợp
\begin{problem}[8DS-HSBN-018]
% nguon: dummy (tự sinh để test)
% lop: 8
% kc_du_kien: 8DS-HSBN-02, 8DS-HSBN-05
% do_kho_du_kien: 4
Tìm tất cả các hàm số bậc nhất $f(x) = ax + b$ ($a \neq 0$) thoả mãn $f\big(f(x)\big) = 4x + 3$ với mọi $x$.

\begin{solution}
Tính $f\big(f(x)\big)$ bằng cách thay $x$ bởi $f(x) = ax + b$ trong công thức của $f$: $f\big(f(x)\big) = a(ax + b) + b = a^2x + ab + b$.

Đẳng thức $a^2x + ab + b = 4x + 3$ phải đúng với mọi $x$. Thay $x = 0$ được $ab + b = 3$ $(1)$; thay $x = 1$ được $a^2 + ab + b = 7$ $(2)$.

Trừ vế theo vế $(2)$ cho $(1)$: $a^2 = 4$; hai số có bình phương bằng $4$ là $2$ và $-2$, nên $a = 2$ hoặc $a = -2$ (đều khác $0$).

Với $a = 2$: từ $(1)$ có $2b + b = 3$, nên $3b = 3$, $b = 1$; ta được $f(x) = 2x + 1$.

Với $a = -2$: từ $(1)$ có $-2b + b = 3$, nên $-b = 3$, $b = -3$; ta được $f(x) = -2x - 3$.

Thử lại với mọi $x$: nếu $f(x) = 2x + 1$ thì $f\big(f(x)\big) = 2(2x + 1) + 1 = 4x + 3$; nếu $f(x) = -2x - 3$ thì $f\big(f(x)\big) = -2(-2x - 3) - 3 = 4x + 6 - 3 = 4x + 3$. Cả hai hàm số đều thoả mãn.

Đáp số: $\boxed{f(x) = 2x + 1 \text{ hoặc } f(x) = -2x - 3}$
\end{solution}
\end{problem}

% ---------------------------------------------------------------- hình bình hành, đường thẳng song song
\begin{problem}[8DS-HSBN-019]
% nguon: dummy (tự sinh để test)
% lop: 8
% kc_du_kien: 8DS-HSBN-04, 8DS-HSBN-05, 8DS-HSBN-06
% do_kho_du_kien: 4
Trên mặt phẳng toạ độ $Oxy$ cho ba điểm $A(0; 1)$, $B(1; 3)$, $C(4; 2)$.

a) Xác định các hàm số bậc nhất có đồ thị lần lượt là đường thẳng $AB$ và đường thẳng $BC$. Chứng minh rằng $A$, $B$, $C$ không thẳng hàng.

b) Tìm toạ độ điểm $D$ sao cho tứ giác $ABCD$ là hình bình hành.

\begin{solution}
a) Đường thẳng $AB$: $y = ax + b$ đi qua $A(0; 1)$ nên $b = 1$; đi qua $B(1; 3)$ nên $3 = a + 1$, suy ra $a = 2$. Vậy $AB: y = 2x + 1$.

a) Đường thẳng $BC$: $y = a'x + b'$ đi qua $B(1; 3)$ và $C(4; 2)$ nên $3 = a' + b'$ và $2 = 4a' + b'$. Từ đẳng thức đầu, $b' = 3 - a'$; thay vào đẳng thức sau: $2 = 4a' + 3 - a' = 3a' + 3$, suy ra $a' = -\dfrac{1}{3}$ và $b' = 3 + \dfrac{1}{3} = \dfrac{10}{3}$. Vậy $BC: y = -\dfrac{1}{3}x + \dfrac{10}{3}$.

a) Thay $x = 4$ vào $y = 2x + 1$ được $y = 9 \neq 2$, nên $C$ không thuộc đường thẳng $AB$; do đó $A$, $B$, $C$ không thẳng hàng.

b) Tứ giác $ABCD$ là hình bình hành khi và chỉ khi các cạnh đối song song: $DC \parallel AB$ và $AD \parallel BC$. Vậy $D$ là giao điểm của đường thẳng $d$ đi qua $C$ song song với $AB$ và đường thẳng $d'$ đi qua $A$ song song với $BC$.

b) $d \parallel AB$ nên $d$ có hệ số góc $2$: $d: y = 2x + c$ với $c \neq 1$; $d$ đi qua $C(4; 2)$ nên $2 = 8 + c$, suy ra $c = -6$ (thoả mãn $c \neq 1$). Vậy $d: y = 2x - 6$.

b) $d' \parallel BC$ nên $d'$ có hệ số góc $-\dfrac{1}{3}$: $d': y = -\dfrac{1}{3}x + e$ với $e \neq \dfrac{10}{3}$; $d'$ đi qua $A(0; 1)$ nên $e = 1$ (thoả mãn). Vậy $d': y = -\dfrac{1}{3}x + 1$.

b) Hai hệ số góc $2 \neq -\dfrac{1}{3}$ nên $d$ và $d'$ cắt nhau; hoành độ giao điểm là nghiệm của $2x - 6 = -\dfrac{1}{3}x + 1$. Nhân hai vế với $3$: $6x - 18 = -x + 3$, nên $7x = 21$, $x = 3$; khi đó $y = 2 \cdot 3 - 6 = 0$. Vậy $D(3; 0)$.

b) Kiểm tra: $D(3; 0)$ thuộc $d'$ vì $-\dfrac{1}{3} \cdot 3 + 1 = 0$; với $D(3; 0)$, các cặp cạnh đối $AB$, $DC$ và $AD$, $BC$ nằm trên các cặp đường thẳng song song, nên $ABCD$ là hình bình hành.

Đáp số: a) $\boxed{AB: y = 2x + 1;\ BC: y = -\dfrac{1}{3}x + \dfrac{10}{3}}$; b) $\boxed{D(3; 0)}$
\end{solution}
\end{problem}

% ---------------------------------------------------------------- khoảng cách lớn nhất từ gốc toạ độ đến đường thẳng
\begin{problem}[8DS-HSBN-020]
% nguon: dummy (tự sinh để test)
% lop: 8
% kc_du_kien: 8DS-HSBN-03, 8DS-HDT-06, 8DS-HDT-01
% do_kho_du_kien: 5
Cho hàm số bậc nhất $y = mx + m + 2$ (với $m \neq 0$) có đồ thị là đường thẳng $d$. Gọi $H$ là chân đường vuông góc kẻ từ gốc toạ độ $O$ đến $d$. Tìm $m$ để độ dài $OH$ (khoảng cách từ $O$ đến $d$) lớn nhất và tính giá trị lớn nhất đó.

\begin{solution}
Trường hợp $m = -2$: $d: y = -2x$ đi qua gốc toạ độ, nên $H$ trùng $O$ và $OH = 0$.

Xét $m \neq 0$ và $m \neq -2$. Cho $x = 0$ được $y = m + 2$, nên $d$ cắt trục tung tại $A(0; m + 2)$ với $OA = |m + 2| > 0$. Cho $y = 0$ được $mx + m + 2 = 0$, suy ra $x = -\dfrac{m + 2}{m}$, nên $d$ cắt trục hoành tại $B\left(-\dfrac{m + 2}{m}; 0\right)$ với $OB = \dfrac{|m + 2|}{|m|} > 0$.

Hai trục toạ độ vuông góc nên tam giác $OAB$ vuông tại $O$; theo định lí Pythagore: $AB^2 = OA^2 + OB^2 = (m + 2)^2 + \dfrac{(m + 2)^2}{m^2} = \dfrac{(m + 2)^2(m^2 + 1)}{m^2}$.

$OH$ là đường cao ứng với cạnh huyền $AB$. Tính diện tích tam giác $OAB$ theo hai cách: $\dfrac{1}{2} \cdot OA \cdot OB = \dfrac{1}{2} \cdot AB \cdot OH$, nên $OH = \dfrac{OA \cdot OB}{AB}$ và $OH^2 = \dfrac{OA^2 \cdot OB^2}{AB^2} = \dfrac{(m + 2)^2 \cdot \dfrac{(m + 2)^2}{m^2}}{\dfrac{(m + 2)^2(m^2 + 1)}{m^2}} = \dfrac{(m + 2)^2}{m^2 + 1}$.

Xét hiệu, dùng hằng đẳng thức bình phương của một tổng, một hiệu: $5 - OH^2 = \dfrac{5(m^2 + 1) - (m^2 + 4m + 4)}{m^2 + 1} = \dfrac{4m^2 - 4m + 1}{m^2 + 1} = \dfrac{(2m - 1)^2}{m^2 + 1}$.

Vì $(2m - 1)^2 \geq 0$ và $m^2 + 1 > 0$ nên $5 - OH^2 \geq 0$, tức là $OH^2 \leq 5$, suy ra $OH \leq \sqrt{5}$. Dấu bằng xảy ra khi và chỉ khi $2m - 1 = 0$, tức là $m = \dfrac{1}{2}$ (thoả mãn $m \neq 0$, $m \neq -2$).

So sánh với trường hợp $m = -2$ ($OH = 0 < \sqrt{5}$), giá trị lớn nhất của $OH$ là $\sqrt{5}$, đạt khi $m = \dfrac{1}{2}$. Thử lại: với $m = \dfrac{1}{2}$, $d: y = \dfrac{1}{2}x + \dfrac{5}{2}$, $OA = \dfrac{5}{2}$, $OB = 5$, $AB^2 = \dfrac{25}{4} + 25 = \dfrac{125}{4}$, nên $OH^2 = \dfrac{OA^2 \cdot OB^2}{AB^2} = \dfrac{625}{4} : \dfrac{125}{4} = 5$ (đúng).

Đáp số: $\boxed{m = \dfrac{1}{2},\ OH_{\max} = \sqrt{5}}$
\end{solution}
\end{problem}

% ---------------------------------------------------------------- vẽ đồ thị y = ax và y = ax + b trên cùng một hệ trục
\begin{problem}[8DS-HSBN-021]
% nguon: dummy (tự sinh để test)
% lop: 8
% kc_du_kien: 8DS-HSBN-03, 8DS-HSBN-04
% do_kho_du_kien: 1
Vẽ trên cùng một mặt phẳng toạ độ $Oxy$ đồ thị $d_1$ của hàm số $y = -\dfrac{1}{2}x$ và đồ thị $d_2$ của hàm số $y = -\dfrac{1}{2}x + 2$ (mô tả cách vẽ bằng lời: nêu hai điểm mà mỗi đường thẳng đi qua). Nêu nhận xét về vị trí của hai đường thẳng $d_1$ và $d_2$.

\begin{solution}
Đồ thị của hàm số bậc nhất $y = ax + b$ ($a \neq 0$) là một đường thẳng, nên để vẽ đồ thị ta chỉ cần xác định hai điểm phân biệt thuộc đồ thị rồi vẽ đường thẳng đi qua hai điểm đó.

Với $d_1\colon y = -\dfrac{1}{2}x$: cho $x = 0$ thì $y = 0$, nên $d_1$ đi qua gốc toạ độ $O(0; 0)$. Vì $d_1$ cắt cả hai trục tại cùng một điểm $O$, ta lấy thêm một điểm khác: cho $x = 2$ thì $y = -\dfrac{1}{2} \cdot 2 = -1$, được điểm $A(2; -1)$. Vậy $d_1$ là đường thẳng $OA$.

Với $d_2\colon y = -\dfrac{1}{2}x + 2$: cho $x = 0$ thì $y = 2$, nên $d_2$ cắt trục tung tại $B(0; 2)$; cho $y = 0$ thì $-\dfrac{1}{2}x + 2 = 0$, suy ra $x = 4$, nên $d_2$ cắt trục hoành tại $C(4; 0)$. Vậy $d_2$ là đường thẳng $BC$.

Hai đường thẳng $d_1$ và $d_2$ có cùng hệ số góc $-\dfrac{1}{2}$ nhưng tung độ gốc khác nhau ($0 \neq 2$), nên $d_1$ song song với $d_2$.

Đáp số: $\boxed{d_1 \text{ đi qua } O(0; 0),\ A(2; -1);\ d_2 \text{ đi qua } B(0; 2),\ C(4; 0);\ d_1 \text{ song song với } d_2}$
\end{solution}
\end{problem}

% ---------------------------------------------------------------- bài toán thực tế: đổi nhiệt độ từ độ C sang độ F
\begin{problem}[8DS-HSBN-022]
% nguon: dummy (tự sinh để test)
% lop: 8
% kc_du_kien: 8DS-HSBN-02, 8DS-PT1-01
% do_kho_du_kien: 2
Số đo nhiệt độ $F$ theo thang Fahrenheit (độ F) và số đo nhiệt độ $C$ theo thang Celsius (độ C) liên hệ với nhau bởi công thức $F = \dfrac{9}{5}C + 32$.

a) $F$ có phải là hàm số bậc nhất của $C$ không? Nếu có, hãy chỉ ra các hệ số $a$, $b$.

b) Tính số đo nhiệt độ theo độ F khi nhiệt độ là $25$ độ C và khi nhiệt độ là $-10$ độ C.

c) Ở nhiệt độ nào thì số đo theo hai thang bằng nhau?

\begin{solution}
a) Công thức $F = \dfrac{9}{5}C + 32$ có dạng $y = ax + b$ (biến $C$ đóng vai trò $x$), với $a = \dfrac{9}{5} \neq 0$ và $b = 32$; vậy $F$ là hàm số bậc nhất của $C$.

b) Thay $C = 25$ vào công thức: $F = \dfrac{9}{5} \cdot 25 + 32 = 45 + 32 = 77$ (độ F).

b) Thay $C = -10$ vào công thức: $F = \dfrac{9}{5} \cdot (-10) + 32 = -18 + 32 = 14$ (độ F).

c) Gọi $t$ là số đo cần tìm, tức là $C = t$ và $F = t$. Thay vào công thức được phương trình $t = \dfrac{9}{5}t + 32$.

c) Chuyển vế: $t - \dfrac{9}{5}t = 32$, tức là $-\dfrac{4}{5}t = 32$, suy ra $t = 32 : \left(-\dfrac{4}{5}\right) = -40$.

c) Thử lại: với $C = -40$ thì $F = \dfrac{9}{5} \cdot (-40) + 32 = -72 + 32 = -40$, đúng bằng $C$.

Đáp số: a) $\boxed{\text{Có; } a = \dfrac{9}{5},\ b = 32}$; b) $\boxed{77 \text{ độ F};\ 14 \text{ độ F}}$; c) $\boxed{-40 \text{ độ (theo cả hai thang)}}$
\end{solution}
\end{problem}

% ---------------------------------------------------------------- bài toán thực tế: ý nghĩa giao điểm của đồ thị với hai trục
\begin{problem}[8DS-HSBN-023]
% nguon: dummy (tự sinh để test)
% lop: 8
% kc_du_kien: 8DS-HSBN-03, 8DS-HSBN-02
% do_kho_du_kien: 2
Lượng xăng $y$ (lít) còn lại trong bình xăng của một ô tô sau $x$ giờ chạy được cho bởi hàm số bậc nhất $y = -6x + 45$ (xe không đổ thêm xăng).

a) Tìm toạ độ giao điểm $A$ của đồ thị hàm số với trục tung và giao điểm $B$ của đồ thị hàm số với trục hoành.

b) Nêu ý nghĩa thực tế của hai điểm $A$, $B$. Xe chạy được nhiều nhất bao lâu thì hết xăng? Trong thực tế, đồ thị của hàm số là phần nào của đường thẳng $AB$?

\begin{solution}
a) Giao điểm với trục tung có hoành độ $x = 0$; thay vào công thức được $y = -6 \cdot 0 + 45 = 45$. Vậy $A(0; 45)$.

a) Giao điểm với trục hoành có tung độ $y = 0$; ta có $-6x + 45 = 0$, suy ra $6x = 45$, $x = \dfrac{15}{2}$. Vậy $B\left(\dfrac{15}{2}; 0\right)$.

b) Điểm $A(0; 45)$ cho biết: lúc xe bắt đầu chạy ($x = 0$), trong bình có $45$ lít xăng.

b) Điểm $B\left(\dfrac{15}{2}; 0\right)$ cho biết: sau $\dfrac{15}{2}$ giờ, tức là $7$ giờ $30$ phút, lượng xăng còn lại bằng $0$. Vậy xe chạy được nhiều nhất $7$ giờ $30$ phút thì hết xăng.

b) Thời gian chạy không âm và lượng xăng còn lại không âm, nên chỉ xét $x$ từ $0$ (lúc xuất phát) đến $\dfrac{15}{2}$ (lúc hết xăng). Do đó trong thực tế đồ thị là đoạn thẳng $AB$ (phần đường thẳng nằm giữa hai trục toạ độ).

Đáp số: a) $\boxed{A(0; 45),\ B\left(\dfrac{15}{2}; 0\right)}$; b) $\boxed{\text{Lúc xuất phát có } 45 \text{ lít; hết xăng sau } 7 \text{ giờ } 30 \text{ phút; đồ thị là đoạn thẳng } AB}$
\end{solution}
\end{problem}

% ---------------------------------------------------------------- tìm tham số khi biết giao điểm với trục tung, trục hoành
\begin{problem}[8DS-HSBN-024]
% nguon: dummy (tự sinh để test)
% lop: 8
% kc_du_kien: 8DS-HSBN-03, 8DS-HSBN-01
% do_kho_du_kien: 2
Cho hàm số bậc nhất $y = (m - 2)x + 2m - 1$ (với $m \neq 2$) có đồ thị là đường thẳng $d$.

a) Tìm $m$ để $d$ cắt trục tung tại điểm có tung độ bằng $5$.

b) Tìm $m$ để $d$ cắt trục hoành tại điểm có hoành độ bằng $-3$.

\begin{solution}
a) Giao điểm của $d$ với trục tung có hoành độ $x = 0$, nên tung độ của nó là $y = (m - 2) \cdot 0 + 2m - 1 = 2m - 1$.

a) Theo đề bài $2m - 1 = 5$, suy ra $2m = 6$, $m = 3$; giá trị này thoả mãn $m \neq 2$. Khi đó $d\colon y = x + 5$ cắt trục tung tại $(0; 5)$.

b) $d$ cắt trục hoành tại điểm có hoành độ $-3$ nghĩa là điểm $(-3; 0)$ thuộc $d$; thay $x = -3$, $y = 0$ vào công thức hàm số: $0 = (m - 2) \cdot (-3) + 2m - 1$.

b) Rút gọn vế phải: $-3m + 6 + 2m - 1 = -m + 5$; do đó $-m + 5 = 0$, suy ra $m = 5$ (thoả mãn $m \neq 2$).

b) Thử lại: với $m = 5$ thì $d\colon y = 3x + 9$; cho $y = 0$ được $3x + 9 = 0$, $x = -3$, đúng như yêu cầu.

Đáp số: a) $\boxed{m = 3}$; b) $\boxed{m = 5}$
\end{solution}
\end{problem}

% ---------------------------------------------------------------- tìm tham số để đường thẳng tạo với hai trục tam giác có diện tích cho trước
\begin{problem}[8DS-HSBN-025]
% nguon: dummy (tự sinh để test)
% lop: 8
% kc_du_kien: 8DS-HSBN-03
% do_kho_du_kien: 3
Cho hàm số bậc nhất $y = (m - 1)x + 3$ (với $m \neq 1$) có đồ thị là đường thẳng $d$. Tìm $m$ để $d$ tạo với hai trục toạ độ một tam giác có diện tích bằng $3$ (đơn vị diện tích).

\begin{solution}
Giao điểm của $d$ với trục tung: cho $x = 0$ thì $y = 3$, được $A(0; 3)$, nên $OA = 3$.

Giao điểm của $d$ với trục hoành: cho $y = 0$ thì $(m - 1)x + 3 = 0$; vì $m \neq 1$ nên $x = -\dfrac{3}{m - 1}$, được $B\left(-\dfrac{3}{m - 1}; 0\right)$, nên $OB = \left|-\dfrac{3}{m - 1}\right| = \dfrac{3}{|m - 1|}$.

Vì $A$, $B$ khác $O$ và lần lượt nằm trên hai trục vuông góc, nên $d$ tạo với hai trục tam giác $OAB$ vuông tại $O$, có diện tích $S = \dfrac{1}{2} \cdot OA \cdot OB = \dfrac{1}{2} \cdot 3 \cdot \dfrac{3}{|m - 1|} = \dfrac{9}{2|m - 1|}$.

Điều kiện $S = 3$ cho $\dfrac{9}{2|m - 1|} = 3$, suy ra $6|m - 1| = 9$, tức là $|m - 1| = \dfrac{3}{2}$.

Hai số có giá trị tuyệt đối bằng $\dfrac{3}{2}$ là $\dfrac{3}{2}$ và $-\dfrac{3}{2}$, nên $m - 1 = \dfrac{3}{2}$ hoặc $m - 1 = -\dfrac{3}{2}$, tức là $m = \dfrac{5}{2}$ hoặc $m = -\dfrac{1}{2}$; cả hai giá trị đều thoả mãn $m \neq 1$.

Thử lại: với $m = \dfrac{5}{2}$, $d\colon y = \dfrac{3}{2}x + 3$ cắt trục hoành tại $(-2; 0)$ và $S = \dfrac{1}{2} \cdot 3 \cdot 2 = 3$; với $m = -\dfrac{1}{2}$, $d\colon y = -\dfrac{3}{2}x + 3$ cắt trục hoành tại $(2; 0)$ và $S = \dfrac{1}{2} \cdot 3 \cdot 2 = 3$.

Đáp số: $\boxed{m = \dfrac{5}{2} \text{ hoặc } m = -\dfrac{1}{2}}$
\end{solution}
\end{problem}

% ---------------------------------------------------------------- điểm thuộc đồ thị cách đều hai trục toạ độ
\begin{problem}[8DS-HSBN-026]
% nguon: dummy (tự sinh để test)
% lop: 8
% kc_du_kien: 8DS-HSBN-01, 8DS-PT1-01
% do_kho_du_kien: 3
Cho hàm số $y = 2x - 6$ có đồ thị là đường thẳng $d$. Tìm toạ độ các điểm $M$ thuộc $d$ sao cho $M$ cách đều hai trục toạ độ.

\begin{solution}
Điểm $M(x_0; y_0)$ thuộc $d$ khi và chỉ khi $y_0 = 2x_0 - 6$, nên $M$ có dạng $M(x_0; 2x_0 - 6)$.

Khoảng cách từ $M(x_0; y_0)$ đến trục hoành bằng $|y_0|$, đến trục tung bằng $|x_0|$. Do đó $M$ cách đều hai trục khi và chỉ khi $|2x_0 - 6| = |x_0|$.

Hai số có giá trị tuyệt đối bằng nhau khi chúng bằng nhau hoặc đối nhau, nên $2x_0 - 6 = x_0$ hoặc $2x_0 - 6 = -x_0$.

Trường hợp $2x_0 - 6 = x_0$: chuyển vế được $x_0 = 6$, khi đó $y_0 = 2 \cdot 6 - 6 = 6$, ta có $M_1(6; 6)$.

Trường hợp $2x_0 - 6 = -x_0$: chuyển vế được $3x_0 = 6$, $x_0 = 2$, khi đó $y_0 = 2 \cdot 2 - 6 = -2$, ta có $M_2(2; -2)$.

Kiểm tra: $M_1(6; 6)$ cách hai trục cùng bằng $6$, $M_2(2; -2)$ cách hai trục cùng bằng $2$, và toạ độ của cả hai điểm đều thoả mãn $y = 2x - 6$.

Đáp số: $\boxed{M_1(6; 6),\ M_2(2; -2)}$
\end{solution}
\end{problem}

% ---------------------------------------------------------------- đưa về dạng y = ax + b, xét vị trí tương đối, tham số bậc hai
\begin{problem}[8DS-HSBN-027]
% nguon: dummy (tự sinh để test)
% lop: 8
% kc_du_kien: 8DS-HSBN-04, 8DS-PTNT-06
% do_kho_du_kien: 3
Cho bốn đường thẳng $d_1\colon y = 2x - 1$; $d_2\colon y = 3x - (x - 3)$; $d_3\colon 4x - 2y - 2 = 0$; $d_4\colon x + 2y + 2 = 0$.

a) Viết các đường thẳng $d_2$, $d_3$, $d_4$ dưới dạng $y = ax + b$ và chỉ ra hệ số góc của mỗi đường thẳng.

b) Xét vị trí tương đối của $d_1$ với mỗi đường thẳng $d_2$, $d_3$, $d_4$.

c) Tìm $m$ để đường thẳng $d\colon y = (m^2 - 2)x + m + 1$ (với $m^2 \neq 2$) song song với $d_1$.

\begin{solution}
a) $d_2$: bỏ ngoặc và thu gọn, $y = 3x - x + 3 = 2x + 3$; hệ số góc bằng $2$.

a) $d_3$: từ $4x - 2y - 2 = 0$ chuyển vế được $2y = 4x - 2$, chia hai vế cho $2$ được $y = 2x - 1$; hệ số góc bằng $2$.

a) $d_4$: từ $x + 2y + 2 = 0$ chuyển vế được $2y = -x - 2$, chia hai vế cho $2$ được $y = -\dfrac{1}{2}x - 1$; hệ số góc bằng $-\dfrac{1}{2}$. Đường thẳng $d_1$ có hệ số góc bằng $2$.

b) Hai đường thẳng $y = ax + b$ và $y = a'x + b'$ song song khi $a = a'$, $b \neq b'$; trùng nhau khi $a = a'$, $b = b'$; cắt nhau khi $a \neq a'$.

b) $d_1$ và $d_2$: $a = a' = 2$ và $b = -1 \neq 3 = b'$, nên $d_1$ song song với $d_2$.

b) $d_1$ và $d_3$: $a = a' = 2$ và $b = b' = -1$, nên $d_1$ trùng với $d_3$.

b) $d_1$ và $d_4$: $2 \neq -\dfrac{1}{2}$ nên $d_1$ cắt $d_4$; hơn nữa hai đường thẳng có cùng tung độ gốc $-1$ nên chúng cắt nhau tại điểm $(0; -1)$ trên trục tung.

c) $d$ song song với $d_1$ khi và chỉ khi $m^2 - 2 = 2$ và $m + 1 \neq -1$.

c) Từ $m^2 - 2 = 2$ suy ra $m^2 - 4 = 0$; phân tích thành nhân tử bằng hằng đẳng thức hiệu hai bình phương: $(m - 2)(m + 2) = 0$, nên $m = 2$ hoặc $m = -2$.

c) Điều kiện $m + 1 \neq -1$ cho $m \neq -2$, nên loại $m = -2$ (khi đó $d\colon y = 2x - 1$ trùng với $d_1$). Với $m = 2$: $m^2 = 4 \neq 2$ và $d\colon y = 2x + 3$ song song với $d_1$.

Đáp số: a) $\boxed{d_2\colon y = 2x + 3;\ d_3\colon y = 2x - 1;\ d_4\colon y = -\dfrac{1}{2}x - 1;\ \text{hệ số góc lần lượt là } 2;\ 2;\ -\dfrac{1}{2}}$; b) $\boxed{d_1 \text{ song song } d_2;\ d_1 \text{ trùng } d_3;\ d_1 \text{ cắt } d_4 \text{ tại } (0; -1)}$; c) $\boxed{m = 2}$
\end{solution}
\end{problem}

% ---------------------------------------------------------------- điểm phụ thuộc tham số luôn chạy trên một đường thẳng cố định
\begin{problem}[8DS-HSBN-028]
% nguon: dummy (tự sinh để test)
% lop: 8
% kc_du_kien: 8DS-HSBN-01, 8DS-HSBN-04
% do_kho_du_kien: 4
Trên mặt phẳng toạ độ $Oxy$, cho điểm $M(m - 1; 2m + 3)$, với $m$ là tham số.

a) Chứng minh rằng khi $m$ thay đổi, điểm $M$ luôn thuộc một đường thẳng $d$ cố định. Viết $d$ dưới dạng $y = ax + b$.

b) Chứng minh rằng không có giá trị nào của $m$ để $M$ thuộc đường thẳng $d'\colon y = 2x + 1$.

c) Tìm $m$ để $M$ thuộc đường thẳng $d''\colon y = -x + 2$.

\begin{solution}
a) Gọi $x_M = m - 1$ và $y_M = 2m + 3$ là toạ độ của $M$. Từ $x_M = m - 1$ suy ra $m = x_M + 1$.

a) Thay $m = x_M + 1$ vào $y_M$: $y_M = 2(x_M + 1) + 3 = 2x_M + 5$. Như vậy với mọi $m$, toạ độ của $M$ thoả mãn đẳng thức $y_M = 2x_M + 5$.

a) Theo điều kiện điểm thuộc đồ thị, $M$ luôn thuộc đồ thị của hàm số $y = 2x + 5$; đường thẳng $d\colon y = 2x + 5$ không phụ thuộc $m$ nên là đường thẳng cố định.

b) $M$ thuộc $d'$ khi và chỉ khi $y_M = 2x_M + 1$, tức là $2m + 3 = 2(m - 1) + 1 = 2m - 1$.

b) Đẳng thức trên tương đương với $3 = -1$, vô lí; vậy không có giá trị nào của $m$ để $M$ thuộc $d'$. Điều này phù hợp với câu a): $d$ và $d'$ có cùng hệ số góc $2$ và tung độ gốc khác nhau ($5 \neq 1$) nên song song, không có điểm chung.

c) $M$ thuộc $d''$ khi và chỉ khi $y_M = -x_M + 2$, tức là $2m + 3 = -(m - 1) + 2 = -m + 3$.

c) Chuyển vế được $3m = 0$, suy ra $m = 0$; khi đó $M(-1; 3)$, và thử lại $-(-1) + 2 = 3$ nên $M$ thuộc $d''$.

Đáp số: a) $\boxed{d\colon y = 2x + 5}$; b) $\boxed{2m + 3 = 2m - 1 \text{ vô lí nên không có } m}$; c) $\boxed{m = 0}$
\end{solution}
\end{problem}

% ---------------------------------------------------------------- xác định hàm số từ tỉ số khoảng cách giao điểm với hai trục
\begin{problem}[8DS-HSBN-029]
% nguon: dummy (tự sinh để test)
% lop: 8
% kc_du_kien: 8DS-HSBN-05, 8DS-HSBN-03, 8DS-HSBN-01
% do_kho_du_kien: 4
Xác định hàm số bậc nhất $y = ax + b$, biết đồ thị của nó đi qua điểm $A(4; 1)$ và cắt trục hoành, trục tung lần lượt tại hai điểm $P$, $Q$ khác gốc toạ độ $O$ sao cho $OP = 2 \cdot OQ$.

\begin{solution}
Đồ thị cắt trục tung tại $Q(0; b)$; vì $Q$ khác $O$ nên $b \neq 0$, và $OQ = |b|$.

Cho $y = 0$ được $ax + b = 0$, $x = -\dfrac{b}{a}$ (vì $a \neq 0$); đồ thị cắt trục hoành tại $P\left(-\dfrac{b}{a}; 0\right)$, nên $OP = \left|-\dfrac{b}{a}\right| = \dfrac{|b|}{|a|}$.

Điều kiện $OP = 2 \cdot OQ$ cho $\dfrac{|b|}{|a|} = 2|b|$; chia hai vế cho $|b| > 0$ được $\dfrac{1}{|a|} = 2$, tức là $|a| = \dfrac{1}{2}$, nên $a = \dfrac{1}{2}$ hoặc $a = -\dfrac{1}{2}$.

Đồ thị đi qua $A(4; 1)$ nên $1 = 4a + b$, suy ra $b = 1 - 4a$.

Với $a = \dfrac{1}{2}$: $b = 1 - 2 = -1 \neq 0$, thoả mãn; hàm số là $y = \dfrac{1}{2}x - 1$, đồ thị cắt trục hoành tại $P(2; 0)$, cắt trục tung tại $Q(0; -1)$, có $OP = 2 = 2 \cdot 1 = 2 \cdot OQ$.

Với $a = -\dfrac{1}{2}$: $b = 1 + 2 = 3 \neq 0$, thoả mãn; hàm số là $y = -\dfrac{1}{2}x + 3$, đồ thị cắt trục hoành tại $P(6; 0)$, cắt trục tung tại $Q(0; 3)$, có $OP = 6 = 2 \cdot 3 = 2 \cdot OQ$.

Đáp số: $\boxed{y = \dfrac{1}{2}x - 1 \text{ hoặc } y = -\dfrac{1}{2}x + 3}$
\end{solution}
\end{problem}

% ---------------------------------------------------------------- tam giác tạo bởi ba đường thẳng, hai cách tính diện tích
\begin{problem}[8DS-HSBN-030]
% nguon: dummy (tự sinh để test)
% lop: 8
% kc_du_kien: 8DS-HSBN-06, 8DS-HSBN-03
% do_kho_du_kien: 4
Cho ba đường thẳng $d_1\colon y = 2x$, $d_2\colon y = -x + 6$ và $d_3\colon y = \dfrac{1}{2}x$ (đơn vị đo trên các trục toạ độ là xentimét).

a) Tìm toạ độ giao điểm $A$ của $d_1$ và $d_2$, giao điểm $B$ của $d_2$ và $d_3$. Chứng minh rằng $d_1$ và $d_3$ cắt nhau tại gốc toạ độ $O$.

b) Tính diện tích tam giác $OAB$.

\begin{solution}[cách 1: hiệu diện tích hai tam giác có đáy trên trục hoành]
a) Hoành độ giao điểm của $d_1$ và $d_2$ là nghiệm của phương trình $2x = -x + 6$; chuyển vế được $3x = 6$, $x = 2$; khi đó $y = 2 \cdot 2 = 4$. Vậy $A(2; 4)$.

a) Hoành độ giao điểm của $d_2$ và $d_3$ là nghiệm của phương trình $\dfrac{1}{2}x = -x + 6$; nhân hai vế với $2$ được $x = -2x + 12$, nên $3x = 12$, $x = 4$; khi đó $y = \dfrac{1}{2} \cdot 4 = 2$. Vậy $B(4; 2)$.

a) Hai đường thẳng $d_1$, $d_3$ có hệ số góc khác nhau ($2 \neq \dfrac{1}{2}$) nên cắt nhau tại đúng một điểm; hoành độ giao điểm là nghiệm của $2x = \dfrac{1}{2}x$, tức là $\dfrac{3}{2}x = 0$, $x = 0$, khi đó $y = 0$. Vậy $d_1$ và $d_3$ cắt nhau tại $O(0; 0)$.

b) Gọi $C$ là giao điểm của $d_2$ với trục hoành: cho $y = 0$ thì $-x + 6 = 0$, $x = 6$, nên $C(6; 0)$ và $OC = 6$ cm.

b) Ba điểm $A$, $B$, $C$ cùng thuộc $d_2$ và hoành độ của $B$ nằm giữa hoành độ của $A$ và của $C$ ($2 < 4 < 6$), nên $B$ nằm giữa $A$ và $C$. Do đó $S_{OAB} = S_{OAC} - S_{OBC}$.

b) Tam giác $OAC$ có đáy $OC = 6$ cm nằm trên trục hoành, chiều cao bằng khoảng cách từ $A$ đến trục hoành là $4$ cm, nên $S_{OAC} = \dfrac{1}{2} \cdot 6 \cdot 4 = 12$ (cm$^2$). Tương tự, tam giác $OBC$ có chiều cao $2$ cm, nên $S_{OBC} = \dfrac{1}{2} \cdot 6 \cdot 2 = 6$ (cm$^2$).

b) Vậy $S_{OAB} = 12 - 6 = 6$ (cm$^2$).

Đáp số: a) $\boxed{A(2; 4),\ B(4; 2),\ d_1 \text{ cắt } d_3 \text{ tại } O(0; 0)}$; b) $\boxed{S_{OAB} = 6 \text{ cm}^2}$
\end{solution}

\begin{solution}[cách 2: bù vào hình vuông]
a) Hoành độ giao điểm của $d_1$ và $d_2$ là nghiệm của phương trình $2x = -x + 6$; chuyển vế được $3x = 6$, $x = 2$; khi đó $y = 2 \cdot 2 = 4$. Vậy $A(2; 4)$.

a) Hoành độ giao điểm của $d_2$ và $d_3$ là nghiệm của phương trình $\dfrac{1}{2}x = -x + 6$; nhân hai vế với $2$ được $x = -2x + 12$, nên $3x = 12$, $x = 4$; khi đó $y = \dfrac{1}{2} \cdot 4 = 2$. Vậy $B(4; 2)$.

a) Hai đường thẳng $d_1$, $d_3$ có hệ số góc khác nhau ($2 \neq \dfrac{1}{2}$) nên cắt nhau tại đúng một điểm; hoành độ giao điểm là nghiệm của $2x = \dfrac{1}{2}x$, tức là $\dfrac{3}{2}x = 0$, $x = 0$, khi đó $y = 0$. Vậy $d_1$ và $d_3$ cắt nhau tại $O(0; 0)$.

b) Xét các điểm $E(4; 0)$, $F(4; 4)$, $G(0; 4)$; tứ giác $OEFG$ là hình vuông cạnh $4$ cm, có diện tích $16$ cm$^2$. Điểm $A(2; 4)$ nằm trên cạnh $GF$ (tung độ bằng $4$, hoành độ từ $0$ đến $4$) và điểm $B(4; 2)$ nằm trên cạnh $EF$ (hoành độ bằng $4$, tung độ từ $0$ đến $4$).

b) Hình vuông được chia thành tam giác $OAB$ và ba tam giác vuông: tam giác $OEB$ vuông tại $E$ có $S = \dfrac{1}{2} \cdot 4 \cdot 2 = 4$; tam giác $BFA$ vuông tại $F$ có $FB = 4 - 2 = 2$, $FA = 4 - 2 = 2$ nên $S = \dfrac{1}{2} \cdot 2 \cdot 2 = 2$; tam giác $AGO$ vuông tại $G$ có $S = \dfrac{1}{2} \cdot 2 \cdot 4 = 4$ (đơn vị cm$^2$).

b) Vậy $S_{OAB} = 16 - (4 + 2 + 4) = 6$ (cm$^2$).

Đáp số: a) $\boxed{A(2; 4),\ B(4; 2),\ d_1 \text{ cắt } d_3 \text{ tại } O(0; 0)}$; b) $\boxed{S_{OAB} = 6 \text{ cm}^2}$
\end{solution}
\end{problem}

% ---------------------------------------------------------------- đường thẳng qua điểm cố định tạo với hai trục tam giác có diện tích nhỏ nhất
\begin{problem}[8DS-HSBN-031]
% nguon: dummy (tự sinh để test)
% lop: 8
% kc_du_kien: 8DS-HSBN-03, 8DS-HSBN-05, 8DS-HDT-01
% do_kho_du_kien: 5
Trên mặt phẳng toạ độ $Oxy$ cho điểm $M(1; 2)$. Xét các đường thẳng $d\colon y = ax + b$ ($a \neq 0$) đi qua $M$, cắt tia $Ox$ tại $A$ và cắt tia $Oy$ tại $B$ ($A$, $B$ khác $O$). Xác định đường thẳng $d$ sao cho tam giác $OAB$ có diện tích nhỏ nhất và tính diện tích nhỏ nhất đó.

\begin{solution}
$d$ đi qua $M(1; 2)$ nên $2 = a \cdot 1 + b$, suy ra $b = 2 - a$.

Giao điểm của $d$ với trục tung là $B(0; b)$; cho $y = 0$ được $x = -\dfrac{b}{a}$, nên giao điểm của $d$ với trục hoành là $A\left(-\dfrac{b}{a}; 0\right)$.

$B$ thuộc tia $Oy$ và khác $O$ nên $b > 0$; $A$ thuộc tia $Ox$ và khác $O$ nên $-\dfrac{b}{a} > 0$, mà $b > 0$ nên $a < 0$. Ngược lại, nếu $a < 0$ thì $b = 2 - a > 0$ và $-\dfrac{b}{a} > 0$; vậy điều kiện của đề bài là $a < 0$.

Đặt $t = -a$, khi đó $t > 0$, $b = 2 + t$, $OB = 2 + t$ và $OA = -\dfrac{b}{a} = \dfrac{2 + t}{t}$.

Tam giác $OAB$ vuông tại $O$ nên diện tích của nó là $S = \dfrac{1}{2} \cdot OA \cdot OB = \dfrac{1}{2} \cdot \dfrac{2 + t}{t} \cdot (2 + t) = \dfrac{(2 + t)^2}{2t}$.

Xét hiệu $S - 4 = \dfrac{(2 + t)^2 - 8t}{2t}$. Khai triển theo hằng đẳng thức bình phương của một tổng: $(2 + t)^2 - 8t = t^2 + 4t + 4 - 8t = t^2 - 4t + 4 = (t - 2)^2$ (hằng đẳng thức bình phương của một hiệu). Vậy $S - 4 = \dfrac{(t - 2)^2}{2t}$.

Vì $(t - 2)^2 \geq 0$ và $2t > 0$ nên $S - 4 \geq 0$, tức là $S \geq 4$; dấu bằng xảy ra khi và chỉ khi $t - 2 = 0$, tức là $t = 2$.

Khi $t = 2$: $a = -2$, $b = 4$, đường thẳng $d\colon y = -2x + 4$ cắt trục hoành tại $A(2; 0)$, cắt trục tung tại $B(0; 4)$, và $S = \dfrac{1}{2} \cdot 2 \cdot 4 = 4$.

Đáp số: $\boxed{d\colon y = -2x + 4,\ S_{\min} = 4 \text{ (đơn vị diện tích)}}$
\end{solution}
\end{problem}

\end{document}
